% Generated draft: requires mathematical and source review.
\documentclass[11pt]{article}
\usepackage[margin=25mm]{geometry}
\usepackage{fontspec}
\setmainfont{Latin Modern Roman}
\usepackage{amsmath,amssymb,mathtools,unicode-math}
\setmathfont{Latin Modern Math}
\usepackage{enumitem,needspace,xurl,hyperref}
\setlength{\parindent}{0pt}
\setlength{\parskip}{5pt}
\setlength{\emergencystretch}{3em}
\newcommand{\N}{\mathbb{N}}
\newcommand{\Z}{\mathbb{Z}}
\newcommand{\Q}{\mathbb{Q}}
\newcommand{\R}{\mathbb{R}}
\newcommand{\C}{\mathbb{C}}
\newcommand{\F}{\mathbb{F}}
\newcommand{\A}{\mathbb{A}}
\newcommand{\PP}{\mathbb{P}}
\newcommand{\E}{\mathbb{E}}
\newcommand{\eps}{\varepsilon}
\newcommand{\dd}{\,\mathrm{d}}
\begin{document}
\section*{UnsolvedMath ID: 10000010}
\textbf{Source problem code:} catalogue entry
\subsection{Definitions and mathematical statement}
For a connected graph $G$ with graph distance $d_G$, one usual $k$-net is a subset $N\subseteq V(G)$ whose distinct points have distance greater than $k$ and such that every vertex is within distance $k$ of $N$. A common net-graph convention connects net points at original distance at most $3k$. The source does not supply its separation, covering, or adjacency constants.

Does there exist a graph admitting unboundedly many successive strict descents through net graphs at suitably chosen scales? In an infinite-chain interpretation, seek graphs $G_0,G_1,\ldots$, scales $k_i$, and nets $N_i$ such that $G_{i+1}$ is the prescribed $k_i$-net graph of $G_i$ and is strictly smaller in the source's intended large-scale sense for every $i$.

The relation strictly smaller is not defined. It cannot silently mean merely non-quasi-isometric under the usual fixed-scale connected-net conventions, since such a net graph gives a coarse model quasi-isometric to its parent. It might concern another size, growth, or net-order relation, but none is specified. The source also does not state the graph class, restrictions on the sequence of scales, or whether unboundedly many means a single infinite chain or arbitrarily long finite chains from one initial graph. These distinctions prevent one unique quantified descent theorem.
\subsection{Short English statement}
Can one graph undergo arbitrarily many strict reductions by repeatedly taking net graphs at chosen scales? The source does not define strict reduction, graph class, or the exact net construction and chain quantifier.
\subsection{Sources}
\begin{itemize}
\item Itai Benjamini, \emph{Coarse Geometry and Randomness}, Saint-Flour notes, Open Problem 1.68, PDF page 15, \url{https://www.wisdom.weizmann.ac.il/~itai/stflouraug24.pdf#page=15}; preserved copy \url{https://arquivo.pt/noFrame/replay/20201231041548id_/http://www.wisdom.weizmann.ac.il/~itai/stflouraug24.pdf}; accessed 2026-07-29. The supplied provenance reports extracted text checked against rendered pages.
\item Local registry \texttt{data/unsolvedmath/problems.json}, permanent UnsolvedMath ID $10000010$, source code catalogue entry, \emph{Unbounded descent through iterated graph nets}. Dataset attribution: Ulam AI and the UnsolvedMath Contributors. The summary cross-refers to AMR-099-0009 for the related nets program, without its statement.
\end{itemize}
\end{document}
